%!TEX TS-program = xelatex
%!TEX encoding = UTF-8 Unicode
%==============================================================================
% Puneeth Prakash - Resume (primary; built by CI)
% Engine: XeLaTeX (required - fontspec)
% Build:  latexmk -xelatex resume.tex   ->  resume.pdf (CI renames it to Puneeth-Prakash-Resume.pdf)
%
% Private build (adds the phone number from the git-ignored private.tex; see
% private.example.tex). Work authorization is NOT private - it is in the
% public document below:
%   latexmk -xelatex -usepretex='\def\privatebuild{}' \
%           -jobname=resume-private resume.tex
%
% Class: awesome-cv.cls (vendored from posquit0/Awesome-CV, LPPL
% v1.3c, with local font-declaration modifications - see the header comment
% in that file for what changed and why).
%
% Fonts: fonts/*.otf (Source Sans 3, SIL OFL 1.1, see
% fonts/LICENSE-SourceSans3.txt). Referenced by explicit file path, not by
% family name, so no fontconfig lookup is required in the CI container.
%==============================================================================
\documentclass[11pt, a4paper]{awesome-cv}

% Emit /ActualText for every character run. Without this, xdvipdfmx builds the
% ToUnicode map from the font's cmap and Source Sans 3 maps U+002D, U+2010 and
% U+2011 to the same glyph, so every hyphen extracted as U+2011 ("cert‑manager")
% and small-caps text extracted with dotless i ("ENGiNEER") - verified with
% pdftotext on the CI-built PDF. With this set, both extract as plain ASCII.
\XeTeXgenerateactualtext=1

\geometry{left=1.1cm, top=0.4cm, right=1.1cm, bottom=0.6cm, footskip=.3cm}

\colorlet{awesome}{awesome-skyblue}
\setbool{acvSectionColorHighlight}{true}
\renewcommand{\acvHeaderSocialSep}{\quad\textbar\quad}

% Tightened from the class defaults (3mm / 2.5mm). Content-top-skip is paired
% with a fixed -3mm/-2mm offset at each call site in the class, so going
% much below ~2mm nets a negative space large enough to pull paragraph
% text up into the section-title line above it - keep this at 2mm or more.
\renewcommand{\acvSectionTopSkip}{1.2mm}
\renewcommand{\acvSectionContentTopSkip}{2mm}

% Header name is 32pt by default; the photo (1.3cm circle) plus that name
% size leaves excess top whitespace. Trim both name lines by 1pt.
\renewcommand*{\headerfirstnamestyle}[1]{{\fontsize{31pt}{1em}\headerfontlight\color{graytext} #1}}
\renewcommand*{\headerlastnamestyle}[1]{{\fontsize{31pt}{1em}\headerfont\bfseries\color{text} #1}}

%-------------------------------------------------------------------------------
% PERSONAL INFORMATION
%-------------------------------------------------------------------------------
% Same headshot used on typedbyme.puneeth.io (assets/images/author/puneeth.png).
\photo[circle,edge,right]{./images/puneeth}
\name{Puneeth}{Prakash}
\position{DevSecOps Engineer{\enskip\cdotp\enskip}Platform Engineer}
% Work authorization rides on the address line, immediately under the job
% title, because in the DACH market it is a go/no-go filter a recruiter looks
% for in the first few seconds. Set upright and in the accent colour so it
% separates from the italic grey address rather than blending into it.
%
% Just the permit name - no "no sponsorship required" or similar gloss. An
% Austrian EU Blue Card is tied to one employer (a new employer still files an
% application), so any such claim would overstate it. The card name alone is
% the fact; the recruiter knows what it means.
\address{Bregenz, Vorarlberg, Austria
  {\enskip\textbullet\enskip}
  {\upshape\bfseries\color{awesome} EU Blue Card (Austria)}}

\email{contact@puneeth.io}
\homepage{typedbyme.puneeth.io}
\github{punitpi}
\linkedin{ppuneeth}

% The phone number must not land in the public, auto-published PDF, so it
% lives in the git-ignored private.tex and is only pulled in when
% \privatebuild is defined on the command line (see header comment).
% Work authorization is public - it is on the address line above.
\ifdefined\privatebuild
  \InputIfFileExists{private.tex}{}{%
    \ClassError{resume}{private build requested but private.tex not found}{Copy private.example.tex to private.tex and fill it in.}}
\fi

%-------------------------------------------------------------------------------
\begin{document}

\makecvheader[C]

\makecvfooter
  {\today}
  {Puneeth Prakash~~~\textperiodcentered~~~R\'esum\'e}
  {\thepage}

%-------------------------------------------------------------------------------
% SUMMARY
%-------------------------------------------------------------------------------
\cvsection{Summary}\par
\begin{cvparagraph}
DevSecOps and platform engineer with 9+ years across cloud infrastructure, Kubernetes, CI/CD, and
backend services on Azure and AWS. Builds the platform layer --- observability, infrastructure as
code, hardened delivery pipelines --- and has shipped the Node.js and .NET services that run on it.
AWS Certified Solutions Architect -- Associate; based in Vorarlberg, Austria.
\end{cvparagraph}

%-------------------------------------------------------------------------------
% SKILLS
%-------------------------------------------------------------------------------
\cvsection{Technical Skills}
\begin{cvskills}
  \cvskill{Cloud \& IaC}{Azure (AKS, Key Vault, App Service, WebJobs), AWS (EKS, ECS Fargate, Lambda, API Gateway, DynamoDB, S3, SQS/SNS), Terraform, CloudFormation}
  \cvskill{Kubernetes \& Observability}{Kubernetes, Helm, Docker, Grafana LGTM stack (Loki, Grafana, Tempo, Mimir), cert-manager, ingress-nginx, Traefik, Inlets Pro}
  \cvskill{DevSecOps \& CI/CD}{GitHub Actions, Jenkins, Harness, Ansible, SonarQube, Black Duck, Keycloak (OAuth2), Tailscale, Authentik}
  \cvskill{Languages \& Backend}{C\# / .NET, JavaScript / Node.js (Express), Python (Flask), Bash, SQL, HCL, React, Microsoft Graph API, MSSQL, MySQL, MongoDB, Redis}
\end{cvskills}

%-------------------------------------------------------------------------------
% EXPERIENCE
%-------------------------------------------------------------------------------
\cvsection{Experience}
\begin{cventries}

  \cventry
    {DevSecOps Engineer / Software Engineer}
    {OMICRON electronics}
    {Klaus, Austria}
    {Dec 2025 -- Present}
    {
      \begin{cvitems}
        \item {Designed and deployed an observability platform on Azure Kubernetes Service (AKS) with the Grafana LGTM stack (Loki, Grafana, Tempo, Mimir), unifying metrics, logs, traces, and alerting.}
        \item {Provisioned the underlying Azure infrastructure from scratch with modular Terraform, making every environment reproducible from code.}
        \item {Packaged the whole stack into a custom umbrella Helm chart for repeatable, versioned deployments.}
        \item {Currently applying DevSecOps practices across the platform: cloud infrastructure security, vulnerability management, and CI/CD pipeline hardening.}
      \end{cvitems}
    }

  \cventry
    {Senior Software Engineer}
    {OneAdvanced}
    {Bengaluru, India}
    {Feb 2021 -- Oct 2025}
    {
      \begin{cvitems}
        \item {Built a secure middleware microservice on AWS ECS Fargate (Node.js, Lambda, DynamoDB) with OAuth2 authentication via Keycloak, handling 100K+ transactions a day at 99.9\% uptime.}
        \item {Deployed and operated production Kubernetes clusters on AWS EKS, running Inlets Pro as a reverse-proxy exit server with cert-manager TLS and zero-touch deployment.}
        \item {Implemented end-to-end CI/CD on GitHub Actions and Harness with infrastructure defined in CloudFormation and Terraform, removing about 80\% of manual release work.}
        \item {Designed reusable .NET Worker Service templates for event-driven backends (logging, encryption, database access) that cut duplicated code by 60\% across projects; built Azure WebJobs webhook endpoints for real-time integration with Microsoft services.}
        \item {Fixed penetration-test and SonarQube findings across services and cut AWS spend by \$1,000+ per month through right-sizing, reserved instances, and automated shutdown of idle resources.}
        \item {Ran internal training sessions on AWS and DevOps practices for the engineering team.}
      \end{cvitems}
    }

  \cventry
    {Software Engineer}
    {OneAdvanced}
    {Bengaluru, India}
    {Mar 2018 -- Feb 2021}
    {
      \begin{cvitems}
        \item {Built a .NET WinForms installer that automated software installation and database updates, cutting setup time from 1 hour to 5 minutes.}
        \item {Co-developed a serverless customer portal on AWS Lambda, API Gateway, and S3, including React front-end components.}
        \item {Rewrote critical SQL stored procedures for a 50\% query-performance gain; automated builds and deployments with Jenkins.}
      \end{cvitems}
    }

  \cventry
    {Test Engineer}
    {Samsung R\&D Institute India}
    {Bengaluru, India}
    {Mar 2017 -- Nov 2017}
    {
      \begin{cvitems}
        \item {Authored and executed 600+ functional, integration, performance, and regression test cases for Project Bixby; managed test-server deployments and analyzed Android crash logs with ADB/DDMS.}
      \end{cvitems}
    }

\end{cventries}

%-------------------------------------------------------------------------------
% PROJECTS
%-------------------------------------------------------------------------------
\cvsection{Projects}
\begin{cventries}

  \cventry
    {Raspberry Pi, Ansible, Docker, Tailscale, Traefik, Authentik, Restic}
    {Homelab Infrastructure \textnormal{\textcolor{graytext}{\small github.com/punitpi/homelab}}}
    {}
    {2024 -- Present}
    {
      \begin{cvitems}
        \item {Four-node, infrastructure-as-code homelab: a cloud edge proxy (Traefik) reaches Raspberry Pi nodes over a Tailscale mesh with no ports exposed to the internet; Authentik SSO, a standby node for failover, and nightly Restic backups to Backblaze B2 with one-command restore.}
        \item {Every service (AdGuard Home DNS, Uptime Kuma monitoring, FreshRSS, code-server, and more) runs as an Ansible-deployed Docker Compose stack, so the whole setup is reproducible and version-controlled.}
      \end{cvitems}
    }

  \cventry
    {Terraform, Azure}
    {Azure Terraform Modules \textnormal{\textcolor{graytext}{\small github.com/punitpi/tf-azure-template}}}
    {}
    {2025}
    {
      \begin{cvitems}
        \item {Reusable, per-environment Terraform modules for resource groups, virtual networks, storage, Key Vault, and App Service, with a remote state backend.}
      \end{cvitems}
    }

  \cventry
    {Hugo, GitHub Actions, LaTeX}
    {TypedByMe Portfolio \textnormal{\textcolor{graytext}{\small typedbyme.puneeth.io}}}
    {}
    {2024 -- Present}
    {
      \begin{cvitems}
        \item {Personal site and technical blog; this résumé is built from LaTeX in CI, published as a GitHub Release, and synced into the site on every push.}
      \end{cvitems}
    }

  \cventry
    {Cloudflare Workers, Node.js, OAuth2}
    {Spotify Now Playing API \textnormal{\textcolor{graytext}{\small github.com/punitpi/spotify-now-playing}}}
    {}
    {2024}
    {
      \begin{cvitems}
        \item {Lightweight proxy for Spotify's Currently Playing API with automatic OAuth token refresh, CORS, and edge caching; deployable to Workers, Docker, or Node.}
      \end{cvitems}
    }

\end{cventries}

%-------------------------------------------------------------------------------
% EDUCATION & CERTIFICATION
%-------------------------------------------------------------------------------
\cvsection{Education}
\begin{cventries}

  \cventry
    {B.E., Computer Science and Engineering}
    {Visvesvaraya Technological University}
    {India}
    {2013 -- 2017}
    {}

\end{cventries}

%-------------------------------------------------------------------------------
% CERTIFICATIONS & AWARDS
%-------------------------------------------------------------------------------
\cvsection{Certifications \& Awards}
\begin{cvhonors}
  \cvhonor{AWS Certified Solutions Architect -- Associate}{Amazon Web Services}{}{2023}
  \cvhonor{Spot Award}{critical project delivered ahead of schedule, OneAdvanced}{}{2023}
\end{cvhonors}

%-------------------------------------------------------------------------------
% LANGUAGES
%-------------------------------------------------------------------------------
% Work authorization is NOT repeated here - it sits on the address line in the
% header, where a recruiter actually looks for it. Stating it twice in one
% two-page document reads as padding.
%-------------------------------------------------------------------------------
\cvsection{Languages}
\begin{cvskills}
  \cvskill{Languages}{English (C1), German (A1, actively learning), Kannada (native)}
\end{cvskills}

%-------------------------------------------------------------------------------
\end{document}
